\chapter*{ABSTRACT}
\addcontentsline{toc}{chapter}{Abstract}

An automated tool can tell you that an application crashed. It cannot tell you
that the application quietly charged the wrong price, because it has no idea
what the right price was. That gap, between detecting a failure and detecting a
\emph{wrong answer}, is why graphical user interface testing is still done by
people, at a cost that grows with every release, and whose hard-won knowledge of an
application's weak points leaves with the tester.

Samsung Research and Development Institute Bangladesh posed the problem and
asked for a demanding change: to move from \emph{directed} testing, where an
agent is given an objective, to \emph{exploratory} testing, where the agent
decides what is worth testing next, infers what correct behaviour would be, and
is better informed next time. The objective stops being an input
and becomes an output, and a failure becomes ambiguous between a real defect, a
limitation of the agent, and an environment that simply blocked it.

The system built in response is an autonomous exploratory testing agent
organised around one persistent knowledge graph. Three stateless
language-model roles (a planner that writes the next test case, an executor that
drives the application, and an investigator that distils each run into
findings) communicate only through that graph. It grows from a requirements
specification, a design export and a defect history, and from the agent's own
observations, so it works where no documentation exists. One planner and one
graph drive both Android and the web. All eight of the partner's requirement
families were delivered, plus ten derived on contact with a real
application. Among them are a fault attribution taxonomy that stops the
agent's own failures being reported as defects, and a safety boundary
around irreversible actions.

Validation combined a regression suite, live campaigns against a production
commerce application and a website, and a reading of seven research
systems. The agent models an application it has never seen, cites the
requirement each test covers, and costs about one US cent per test, orders of
magnitude below published figures. Autonomy and requirement coverage
remain modest and are reported unadjusted. The review also yielded a finding
that constrains every claim in this field, this project's included: none of the
seven systems reports precision or recall for defect detection, and the two that
do injected their faults by hand, because without a specification there is no
oracle beyond a crash, and without an oracle no ground truth. Grounding the
agent in a specification attacks that problem directly; the seeded-defect
experiment that would close it is the most valuable work remaining.

\vspace{0pt}
\noindent\textbf{Keywords:} exploratory testing; GUI test automation; LLM agents; knowledge graphs.
\clearpage
